\documentclass[10pt,a4paper]{amsart}
\usepackage[margin=19mm]{geometry}
\usepackage{amsmath,amssymb,mathtools}
\usepackage[scr=rsfs,cal=euler]{mathalfa}
\usepackage{xfrac}
\usepackage[unicode,hidelinks,bookmarks=false]{hyperref}
\newcommand{\ie}{i.e.\ }
\newcommand{\cf}{cf.\ }
\newcommand{\resp}{respectively\ }
\newcommand{\Subsection}{\subsection}
\providecommand{\nobreakdash}{\nobreak\hbox{-}}
\setlength{\emergencystretch}{2em}
\multlinegap0pt
\usepackage{siunitx}
\usepackage{siunitx}
\usepackage{siunitx}
\usepackage{siunitx}
\usepackage{tikz}
\usepackage{amssymb}
\usepackage{siunitx}
\newtheorem{proposition}{Proposition}
\title{The Parabolic Mellin Transform: Gamma and Zeta Integral Representations}
\author{Peter Reinhard Hansen$^{\mathsection}$ \and Chen Tong$^{\ddagger}$}
\date{Source: arXiv:2602.17007v2 (CC BY 4.0)}
\begin{document}
\maketitle

\noindent\emph{Source and licence.} The Parabolic Mellin Transform: Gamma and Zeta Integral Representations. Version arXiv:2602.17007v2 (CC BY 4.0), distributed by arXiv under the Creative Commons Attribution 4.0 International licence, http://creativecommons.org/licenses/by/4.0/, as recorded in the arXiv licence metadata for this exact version and shown on the abstract page https://arxiv.org/abs/2602.17007v2. Full licence notice: \texttt{licenses/arxiv-2602.17007-CC-BY-4.0.txt}.\par\smallskip

\noindent\textit{Editorial scope.} Counted unit: the article's proposition prop:X-sufficient-RH (source0.tex lines 815--820). The article's complete original proof of it is delivered with it (source0.tex lines 822--849, one \texttt{proof} environment of 28 lines). No mathematics is written, repaired or re-derived: the statement and the proof are the article's own lines, in the article's own notation, and no retained line is substituted or re-flowed.  No further source-local context is needed: every cross-reference of every retained line resolves inside the two delivered spans.  Nothing was omitted from any retained span, so the package carries no omission record.  No retained line contains a citation, so the wrapper carries no bibliography and no bibliography entry.  No retained line contains a figure environment, an image-inclusion command or drawing code, so no image is delivered and the package declares no asset dependency.  Editorial formatting only, in addition: an AMS \texttt{amsart} preamble that restates the article's own declarations and macros used by the retained lines, namely the environments \texttt{proposition} on separate counters, together with the standard packages those lines need.  The retained environments print in this document as Proposition 1 in order of appearance, whereas the article prints the same items with the numbering of its own full document.  The complete native source of the article as distributed in the arXiv e-print for this exact version is bundled unchanged as \texttt{sources/194936/source0.tex}.
\begin{proposition}[Auxiliary sufficient condition for RH]
\label{prop:X-sufficient-RH}
If every zero of $\mathcal{X}(\tau)$ in the strip
$|\operatorname{Im}(\tau)|<\tfrac{1}{2}$ is real, then the Riemann hypothesis
holds.
\end{proposition}

\begin{proof}
Using
$$
\frac{1}{e^{-u}-1}+\frac{1}{e^{-u}+1}
=
-\frac{1}{\sinh(u)},
$$
we obtain, with $u=u_t$,
$$
\mathcal{S}(z)
=
-\frac{1}{2}
\int_{-\infty}^{\infty}
\frac{u_t^z-u_t^{-z}}{\sinh(u_t)}dt
=
-\int_{-\infty}^{\infty}
\frac{\sinh(z\log u_t)}{\sinh(u_t)}dt.
$$
Therefore $\mathcal{S}(i\tau)=-i\mathcal{X}(\tau)$.

Let $\rho=\tfrac{1}{2}+z$ be a nontrivial zero of $\zeta$, so
$|\operatorname{Re}(z)|<\tfrac{1}{2}$. By the preceding paragraph,
$\mathcal{S}(z)=0$. Writing $z=i\tau$, equivalently $\tau=-iz$, gives
$|\operatorname{Im}(\tau)|<\tfrac{1}{2}$ and
$\mathcal{X}(\tau)=0$ since $\mathcal{S}(i\tau)=-i\mathcal{X}(\tau)$. By hypothesis, this zero
$\tau$ is real. Hence $z=i\tau$ is purely imaginary, so
$\operatorname{Re}(\rho)=\tfrac{1}{2}$. This proves the Riemann hypothesis (RH).
\end{proof}

\end{document}
